% =============================================================================
% GUIDA ALLA SCRITTURA DELLA RICERCA - SOLO PER IL MODELLO - ITALIANO
% Disabilitarla in configuration.tex con:
% \providecommand{\ISISResearchWritingGuideChoice}{disabled}
% Il capitolo adatta indicazioni generali per articoli scientifici a una tesi
% più estesa. E materiale didattico e non deve apparire nella tesi finale.
% =============================================================================

\chapter{Guida del modello: scrivere una tesi di ricerca efficace}\label{ch:writing-guide-it}

\begin{keypoint}[Capitolo presente solo nel modello]
Questo capitolo mostra come trasformare un'idea di ricerca in una tesi leggibile.
Non appartiene al progetto scientifico di esempio. Prima della consegna,
disabilitarlo in \texttt{configuration.tex} impostando
\texttt{\textbackslash ISISResearchWritingGuideChoice} su \texttt{disabled}.
\end{keypoint}

Le indicazioni adattano alla tesi la presentazione di Simon Peyton Jones per
Microsoft Research, \emph{How to Write a Great Research Paper}. La fonte e
consultabile sul sito di
\href{https://www.microsoft.com/en-us/research/wp-content/uploads/2016/07/How-to-write-a-great-research-paper.pdf}{Microsoft Research}.
Una tesi è più lunga di un articolo da conferenza e può includere diversi studi,
ma beneficia comunque di una tesi centrale, una narrazione esplicita,
contributi verificabili, esempi concreti ed evidenze collegate a ogni
affermazione importante.

\section{Partire da una tesi centrale}
Una tesi efficace può normalmente essere riassunta in una frase che dichiara
che cosa è stato appreso, costruito, dimostrato o spiegato. Piu contributi
possono sostenere la frase, ma non devono apparire come progetti scollegati.

\begin{researchquestion}[Tesi centrale -- esempio]
Un'architettura Edge-AI modulare può ridurre il tempo di risposta e l'uso delle
risorse, mantenendo prestazioni predittive utili, se acquisizione, inferenza,
politica di sicurezza e presentazione vengono valutate come componenti separati.
\end{researchquestion}

\begin{keypoint}[Adattamento dall'articolo alla tesi]
Per un articolo può bastare un'unica idea molto netta. Per una tesi, mantenere
una tesi centrale e organizzare attorno ad essa da due a quattro contributi
secondari. Ogni contributo deve essere abbastanza specifico da consentire al
lettore di verificare se e stato effettivamente realizzato.
\end{keypoint}

Una possibile narrazione per una tesi tecnica è:
\begin{enumerate}
  \item un problema concreto che il lettore possa immaginare;
  \item la ragione per cui il problema conta e il requisito ancora insoddisfatto;
  \item l'idea della tesi e le ipotesi entro cui si applica;
  \item il metodo o il sistema che realizza l'idea;
  \item le evidenze che sostengono o limitano l'affermazione;
  \item un confronto corretto con gli approcci precedenti;
  \item conclusioni, limiti e domande aperte.
\end{enumerate}

\section{Pagina di esempio: un sommario in quattro frasi}
Il sommario deve permettere di ricostruire problema, approccio, evidenza
principale e contributo senza aprire il resto della tesi. L'esempio seguente usa
quattro frasi; i valori numerici sono segnaposto.

\begin{frontmatterbox}
\textbf{Problema.} I dispositivi edge negli ambienti intelligenti devono
rispondere rapidamente, pur operando con memoria limitata e con un confine di
sicurezza chiaramente definito. \textbf{Approccio.} Questa tesi progetta un
prototipo modulare che separa acquisizione, inferenza locale, applicazione delle
politiche e presentazione, e valuta ogni componente con un protocollo
riproducibile. \textbf{Evidenza.} Nella valutazione di esempio, la configurazione
ottimizzata raggiunge il 91.0\% di accuratezza con latenza mediana di 29~ms e
241~MB di memoria, rispetto a 72.0\%, 84~ms e 410~MB della baseline.
\textbf{Contributo.} I risultati mostrano come un'architettura collegata alle
evidenze renda espliciti i compromessi tra prestazioni, risorse e sicurezza;
tutti i valori devono essere sostituiti con i risultati reali dello studente.
\end{frontmatterbox}

\begin{keypoint}[Test del sommario]
Dopo aver letto solo il sommario, il relatore dovrebbe poter sottolineare un
problema, un metodo, un risultato principale e un contributo. Eliminare frasi
generiche come ``la tecnologia e sempre più importante'' quando non definiscono
direttamente il problema di ricerca.
\end{keypoint}

\section{Pagina di esempio: introdurre il problema con un caso guida}
Iniziare con un caso piccolo invece che con una rassegna molto ampia. Il caso
dona uno scopo ai termini tecnici e aiuta il lettore a ricordare la tensione
centrale.

\begin{resultbox}[Esempio di paragrafo iniziale]
Alle 02:14, una telecamera in un locale industriale rileva un movimento vicino a
un armadio riservato. Il collegamento di rete è temporaneamente assente, quindi
un sistema esclusivamente cloud non può restituire una decisione entro il
vincolo di 50~ms. Un modello locale può rispondere in tempo, ma solo se usa meno
di 300~MB e mantiene il tasso di falsi allarmi entro il limite operativo. La
tesi chiede se una pipeline edge modulare possa soddisfare insieme i tre
vincoli.
\end{resultbox}

Il paragrafo svolge quattro funzioni: presenta una situazione concreta, rende
visibile il vincolo, identifica la lacuna e conduce a una domanda di ricerca.
Non apre con molte pagine di affermazioni generiche sull'importanza
dell'intelligenza artificiale.

\section{Pagina di esempio: dichiarare contributi verificabili}
Scrivere presto l'elenco dei contributi e usarlo per decidere che cosa appartiene
alla tesi. Evitare contributi che si limitano a dire che un sistema viene
descritto o che un tema viene studiato. Preferire affermazioni che indicano un
artefatto, una proprietà, un confronto o un risultato misurato.

\begin{keypoint}[Esempio di elenco dei contributi]
\begin{enumerate}
  \item Definiamo un requisito misurabile che combina latenza, memoria e vincoli
        di sicurezza, documentandone le ipotesi nel capitolo di formulazione del
        problema.
  \item Implementiamo un'architettura a quattro stadi i cui componenti possono
        essere verificati separatamente, fornendo interfacce e configurazioni
        necessarie alla ricostruzione.
  \item Confrontiamo il prototipo con una baseline dichiarata mediante prove
        ripetute e riportiamo accuratezza, latenza, memoria, incertezza e casi di
        fallimento.
  \item Identifichiamo le condizioni operative in cui l'approccio non raggiunge
        il requisito, invece di presentare soltanto la configurazione migliore.
\end{enumerate}
\end{keypoint}

\subsection{Collegare ogni affermazione a un'evidenza}
L'introduzione formula promesse; il corpo della tesi deve fornire evidenze per
ciascuna promessa. Una matrice affermazione-evidenza è uno strumento semplice di
pianificazione.

\begin{table}[htbp]
\centering
\caption{Esempio di mappa affermazione-evidenza per una tesi.}
\label{tab:guide-claim-evidence-it}
\begin{tabularx}{\textwidth}{p{0.23\textwidth}Yp{0.27\textwidth}}
\toprule
Affermazione & Evidenza richiesta & Collocazione prevista\\
\midrule
Il prototipo rispetta il vincolo di latenza & Misure ripetute, distribuzione e
confronto con la baseline & Protocollo, figura sulla latenza e risultati grezzi
in appendice\\
L'architettura e modulare & Definizione delle interfacce, test dei componenti ed
esperimento di sostituzione & Capitolo di progettazione e test a livello di
componente\\
Il metodo è robusto a un ingresso mancante & Condizioni di corruzione definite
prima dei test, perdita di prestazione e incertezza & Esperimento di robustezza e
analisi dei fallimenti\\
Il risultato è riproducibile & Dati, codice, parametri, seed, hardware e comando
di riproduzione versionati & Metodologia, appendice e checklist facoltativa\\
\bottomrule
\end{tabularx}
\end{table}

Usare riferimenti in avanti da ogni contributo alla sezione, tabella,
dimostrazione, caso di studio o esperimento che lo sostiene. Un paragrafo
meccanico che elenca l'ordine dei capitoli è meno utile di una narrazione che
collega affermazioni ed evidenze.

\section{Pagina di esempio: intuizione prima dei dettagli formali}
Spiegare l'idea come se venisse disegnata su una lavagna. Introdurre la notazione
solo dopo che il lettore ha compreso il ruolo dei componenti e la ragione per
cui sono necessari.

\begin{figure}[htbp]
\centering
\pipelinebox{Osserva la\\qualità}
\pipelinearrow
\pipelinebox{Stima la\\affidabilità}
\pipelinearrow
\pipelinebox{Pesa le\\evidenze}
\pipelinearrow
\pipelinebox{Predice e\\dichiara i limiti}
\caption{Spiegazione a livello di lavagna prima di equazioni e dettagli implementativi.}
\label{fig:guide-intuition-it}
\end{figure}

\begin{keypoint}[Prima e dopo]
\textbf{Troppo presto:} ``Sia $z_m=f_m(x_m)$ e si ottimizzi un obiettivo
multimodale regolarizzato.'' La notazione arriva prima del problema.\par
\textbf{Meglio:} ``Ogni sensore propone una risposta e indica quanto e affidabile
il segnale corrente. Il modulo di fusione riduce il peso di un sensore rumoroso
o assente. Formalizziamo questo comportamento dopo il caso guida.''
\end{keypoint}

Quando l'intuizione è stabile, definire simboli, ipotesi, algoritmi e
dimostrazioni con precisione. Gli esempi non sostituiscono il rigore; preparano
il lettore a comprenderlo.

\section{Usare frasi attive e dirette}
Preferire frasi che dichiarano chi compie l'azione. Per esempio, scrivere
``Abbiamo eseguito 30 prove indipendenti'' invece di ``Sono state eseguite 30
prove'', e ``Il modello ottimizzato ha usato 241~MB'' invece di ``E stata
osservata una riduzione nell'utilizzo della memoria''. Usare la forma passiva
quando l'attore e davvero irrilevante, non come stile predefinito.

\clearpage
\section{Pagina di esempio: collocare lo stato dell'arte dove aiuta}
Una tesi richiede spesso più contesto di un articolo e il corso di laurea può
imporre una rassegna iniziale. All'inizio mantenere soltanto i concetti necessari
per comprendere il problema. Collocare il confronto dettagliato con le
alternative dopo che il lettore ha capito l'idea proposta, oppure ritornare
sulla letteratura in un capitolo successivo di posizionamento.

\begin{table}[htbp]
\centering
\caption{Esempio di sintesi: confrontare le ipotesi invece di elencare articoli.}
\label{tab:guide-related-it}
\begin{tabularx}{\textwidth}{Ycccc}
\toprule
Approccio & Decisione locale & Ingresso mancante & Calibrazione & Risorse\\
\midrule
Baseline solo cloud & No & Servizio interrotto & Facoltativa & Solo server\\
Fusione anticipata & Sì & Limitata & Rara & Talvolta\\
Fusione tardiva & Sì & Moderata & Facoltativa & Spesso\\
Approccio della tesi & Sì & Esplicito & Richiesta & Device e server\\
\bottomrule
\end{tabularx}
\end{table}

Descrivere i lavori precedenti in modo generoso e accurato. Dichiarare che cosa
un approccio realizza bene, le ipotesi entro cui funziona e la dimensione precisa
in cui la tesi lo estende o se ne differenzia. Riconoscere i limiti del nuovo
approccio rende il confronto più credibile.

\section{Pagina di esempio: riportare evidenze, non soltanto il valore migliore}
Il capitolo dei risultati deve separare osservazione e interpretazione.

\begin{resultbox}[Osservazione]
Su cinque seed illustrativi, la latenza mediana diminuisce da 84~ms a 29~ms.
L'intervallo interquartile diminuisce da 12~ms a 5~ms. Una configurazione
hardware supera il limite quando il throttling termico inizia dopo 20 minuti.
\end{resultbox}

\begin{keypoint}[Interpretazione]
La mediana più bassa e la dispersione ridotta sostengono l'affermazione sulla
reattività nell'hardware testato. Il fallimento dovuto al throttling limita
l'affermazione ai carichi e alle condizioni di raffreddamento rappresentati dal
protocollo. La tesi non deve estendere il risultato a dispositivi non testati.
\end{keypoint}

Riportare incertezza, risultati negativi, deviazioni dal protocollo e
spiegazioni alternative. Distinguere le analisi esplorative dai test pianificati
in anticipo. Un limite non è una scusa: definisce il confine del risultato.

\clearpage
\section{Pagina di revisione: verificare il percorso del lettore}
Il lettore dovrebbe poter seguire la tesi centrale leggendo titolo, sommario,
introduzione, aperture dei capitoli, didascalie, sintesi dei risultati e
conclusioni. Eseguire la seguente revisione prima di perfezionare la tipografia.

\begin{enumerate}
  \item Scrivere la tesi centrale in una sola frase.
  \item Verificare che il sommario dichiari problema, approccio, risultato e contributo.
  \item Sostituire un'apertura troppo ampia con un caso guida concreto.
  \item Rendere ogni contributo verificabile e indicare la relativa evidenza.
  \item Spiegare l'intuizione prima di notazione, algoritmi e dettagli implementativi.
  \item Spostare il confronto dettagliato con lo stato dell'arte dopo l'idea della
        tesi, salvo diversi requisiti istituzionali.
  \item Separare osservazioni misurate, interpretazione e speculazione.
  \item Riportare incertezza, limiti, risultati negativi e condizioni di validità.
  \item Preferire linguaggio attivo e diretto e figure informative.
  \item Chiedere a un lettore esperto e a uno non esperto dove si sono persi,
        quindi correggere il testo invece di attribuire il problema al lettore.
\end{enumerate}

\begin{keypoint}[Verifica finale specifica per la tesi]
La presentazione di Microsoft Research riguarda articoli scientifici, non i
regolamenti universitari. Conservare capitoli obbligatori, dichiarazioni,
formattazione, requisiti etici e criteri di valutazione del corso di laurea.
Applicare i principi narrativi entro tali vincoli istituzionali.
\end{keypoint}
